\section{Friendship, authority, and controversy}\label{sec:jerome-controversies}

\subsection{Rufinus and the Origenist rupture}

Rufinus of Aquileia belonged to Jerome's early ascetic friendship. Both men
traveled east, learned from Greek Christianity, translated Origen, and joined
Latin ascetic networks sustained by elite women. Rufinus lived in Jerusalem
with Melania the Elder while Jerome settled at Bethlehem with Paula. Their
later war was therefore not a clash between an Origenist stranger and a man
untouched by Origen. It was a struggle over how to receive a formative teacher,
which propositions to reject, who had concealed or corrected them, and which
translator could control the record of his own debt.

The Palestinian Origenist controversy intensified in the 390s. Epiphanius of
Salamis condemned teachings attributed to Origen; John, bishop of Jerusalem,
resisted the campaign and clashed with Jerome's circle. Epiphanius's ordination
of Jerome's brother Paulinian outside John's ordinary control added a concrete
jurisdictional grievance. Jerome's \emph{Against John of Jerusalem}, Letter 84,
and later apologies defend his orthodoxy. Their detail is evidence for the
conflict, but not a neutral summary of John's position.

Rufinus's Latin rendering of Origen's \emph{On First Principles} in 398 and its
preface associated Jerome with a similarly selective translation practice.
Jerome answered with a more literal rival rendering and public denunciation.
Jerome's three books \emph{Against Rufinus} preserve charges,
countercharges, private letters turned into exhibits, and learned insult in
Jerome's arrangement. Rufinus's \emph{Apology against Jerome} 2 now supplies a
controlled counter-archive at decisive points. Rufinus makes Jerome's own
Letter 22 an exhibit: he attacks its satire of clergy and ascetics, charges
that the Ciceronian dream's renunciation conflicts with Jerome's later teaching
of classical authors, and contests the Hebrew-based revision as an injury to
the Church's received Greek Scriptures. These are an opponent's arguments, not
neutral verdicts; they prove that Jerome's self-revision and textual program
were disputed in precise terms rather than by vague personal jealousy. When
the Alexandrian bishop Theophilus moved against Origenist monks
around 399--400, Jerome celebrated an alliance that further isolated his
former friends.

Some disputed propositions were serious: the pre-existence and restoration of
souls, the resurrection body, and the final fate of the devil bore on Christian
doctrine. Yet doctrinal gravity does not sanctify every tactic. Jerome reshaped
the story of his own dependence, attributed motives he could not inspect, and
used ridicule as a weapon. Historical judgment requires both archives and
refuses the victor's retrospective simplicity. Book 2 of Rufinus is now
controlled, while the full reciprocal corpus and institutional sequence still
limit judgment.

\subsection{John of Jerusalem and the limits of local power}

Jerome's international reputation could not erase ecclesial geography.
Bethlehem lay within the jurisdiction of Jerusalem's bishop. Jerome was a
presbyter, not the autonomous bishop of a scholarly republic. Disputes over
communion, Paulinian's ordination, buildings, and access to worship exposed the
tension between charismatic textual authority and local episcopal government.
Jerome's appeals to Rome and powerful friends made the conflict Mediterranean.
They did not make jurisdiction irrelevant.

The episode also complicates later claims that Jerome consistently stood above
party. He could defend Roman communion in the Antiochene schism, support
Epiphanius against John, and later use Theophilus of Alexandria as an orthodox
ally. Each alignment had theological reasons and institutional consequences.
``The Church condemned Origen'' is too undifferentiated to explain who acted,
when, against which propositions, and through what local authority.

\subsection{Augustine: a friendship entirely in letters}

Augustine and Jerome never met. Their relationship was made by copied letters,
messengers, delays, and accidental publication. Around 394 Augustine challenged
Jerome's interpretation of Galatians 2:11--14. Jerome followed a traditional
reading in which Peter and Paul's conflict was staged for pastoral effect;
Augustine feared that attributing a useful falsehood to an apostle would weaken
Scriptural truth. A letter intended as private correction circulated before
Jerome received it, making intellectual disagreement a question of honor.

Their long Letter 82/116 exchange does not dissolve the disagreement. Augustine
argues that Paul's rebuke was truthful and asks for frank correction between
friends; Jerome answers with erudition, rank-conscious courtesy, and eventual
desire for peace. The so-called gourd dispute over Jonah's plant likewise
joined Hebrew philology to a congregation's received wording. Augustine did
not reject Hebrew study; he feared pastoral disruption and valued the
Septuagint's ecclesial authority.

The relationship matured. They exchanged questions about the soul and later
opposed Pelagian positions, but even alliance did not make their methods
identical. Their correspondence is an unusually valuable N as well as A
dossier: each author controls his own letter, while the exchange exposes the
other's claims to challenge.

\subsection{Marriage, virginity, and the price of victory}

Jerome's \emph{Against Jovinian} (393) opposed claims he understood to level
the reward of virginity and marriage, diminish fasting, and overstate baptismal
security. He defended voluntary virginity and ascetic discipline within a
Christian tradition larger than himself. The controversy also displays his
habit of making hierarchy sound like contempt. Even allies worried that his
comparisons injured Christian marriage; Letter 49 acknowledges Pammachius's
attempt to limit circulation while defending the work.

\emph{Against Helvidius} defended Mary's perpetual virginity and the ascetic
honor of virginity. Its doctrinal conclusion became Catholic reception; its
philological arguments and invective remain the work of a fourth-century
controversialist. \emph{Against Vigilantius} (406) defended veneration of
martyrs' relics, vigils, almsgiving to Jerusalem, and celibate asceticism. It
also attacks Vigilantius's occupation, region, appearance, and alleged drinking
with language so abusive that the opponent's actual pastoral concerns can
disappear. The treatise must be read as evidence for developing cult and
ascetic practice and for Jerome's failure of proportion.

Women gained genuine alternatives through ascetic communities: education,
property redirected to charity, leadership outside remarriage, and sustained
scriptural work. At the same time, Jerome's praise could make a woman's body a
symbol in male controversy, treat family affection as weakness, and presume
authority over wealth he did not own. Letter 130 to Demetrias, written thirty
years after Letter 22, is more moderated about fasting and work. Development
matters, but it does not delete earlier words or their effects.

\subsection{Pelagian controversy and physical danger}

Around 415--417 Jerome composed the \emph{Dialogue against the Pelagians}, presenting
the positions through literary interlocutors rather than a neutral transcript
of Pelagius. He emphasized the necessity of grace and the possibility of
post-baptismal sin. His alliance with Augustine now mattered more than their
old friction.

Augustine, Innocent, and Jerome report that the Bethlehem monastic complex was
attacked and burned in 416. Augustine uses report language while alleging a
deacon's death (\emph{Proceedings of Pelagius} 66). Innocent's Letters 136--137
express sympathy and blame local episcopal failure, but also say that no
individual accusation had been made and that reports remained rumor rather
than a pleaded case. Jerome's Letters 138--139 describe danger, displacement,
and the destruction of his house's material resources in openly partisan
language. The convergence supports serious violence and physical destruction;
it does not establish a casualty register, perpetrator, legal finding, command
responsibility, or rebuilding history.

\begin{studyframe}
\textbf{Controversy as biography.} Jerome's disputes were not detachable bad
manners around a pure scholarly core. They shaped which books he translated,
which alliances sustained the workshop, which opponents survived in Latin,
and how he narrated himself. Zeal for truth explains the stakes; it does not
make humiliation, caricature, or uncertain accusation into virtues.
\end{studyframe}
